\section{Comparison with the tiling-specific automata}
\label{sec:comparison}

It is worth setting this construction beside Margenstern's exactly, since the two answer different questions and
the difference is the point.

\subsection{One automaton against a family}

The tiling-specific automata are a family. There is a pentagrid automaton, a heptagrid automaton, a dodecagrid
automaton, each a separate object with its own rules \cite{margenstern2014penta5, margenstern2009hepta,
margenstern2021dodeca5}. Adding a new tiling to the family means designing a new automaton and checking a new
table. The present construction is a single automaton. Adding a new tiling means nothing at all, since the
rules already apply to its cell graph. Where the family covers the tilings it has reached, the single automaton
covers every regular tiling, including all the ones never individually reached.

\subsection{Rotation invariance against oriented switches}

The deeper difference is how each construction recovers direction without a global frame. Margenstern's
automata are rotation invariant. A rule holds under every rotation of the cell, so a cell need not know its
orientation, and direction is encoded in the states themselves, with the help of milestones, small markings
that decorate a track and pick out its sense locally. Rotation invariance is a strong and elegant property, and
it is the reason those automata read as genuine cellular automata in the strict sense. The cost is that the
rotations of a cell, and how they permute its neighbours, are different in each tiling, so rotation invariance
must be re-established for every tiling, which is much of why each tiling needs its own table.

The present automaton is not rotation invariant. Its switch cells are oriented, carrying labels for the trunk
and the two branches as part of the initial configuration. This is the price of uniformity. Because the labels
are local data attached to the configuration, not a global frame, the automaton is still tiling-independent,
but it is no longer rotation invariant. The status of this is the same as that of any weakly universal
automaton with a structured initial configuration. The track network is the program, and the labels are part
of the program, laid down once and never changing.

We regard the two as complementary. If one wants the smallest rotation-invariant automaton on a fixed tiling,
the tiling-specific constructions are the answer, and they have been pushed to remarkably few states, two on
the pentagrid \cite{margenstern2015penta2}, four on the heptagrid \cite{margenstern2009hepta}, four on the
dodecagrid \cite{margenstern2021dodeca4}. If one wants a single object that is universal on the whole family at
once, the present construction is the answer, at the cost of oriented switches.

\subsection{The locomotive's direction trick}

One small point of technique deserves note, because it is what lets the dynamic alphabet stay so small while
remaining tiling-independent. The head needs to move forward and not backward, and on a bare graph there is
nothing to tell it which way is forward. The two-cell locomotive supplies the answer at no extra cost. The cell
behind the head is the trailing cell, which is not clear track, so the only clear track cell adjacent to the
head is the one ahead, by Lemma~\ref{lem:direction}. Forward is read off the local pattern rather than off the
plane. This is the same service that Margenstern's milestones provide, achieved here by the shape of the
locomotive itself, and it is the reason the construction needs no decoration of the track and no orientation of
plain track cells.
